\cotefile{100}
\scanpage{41}
\uline{Remarque} \deleted{\unsure{Supp.|Supposons}} Soit \deleted{\unsure{B}}$S$ \inserted{$=\unsure{\mathrm{Spm}|\mathrm{Spf}|\mathrm{Spec}}\,R$} un schéma local noeth. complet,
$\mathfrak X$ un schéma formel propre \deleted{et} plat sur $S$,
tel que $H^0(X_0,\mathcal O_{X_0})\simeq k$ (pour simplifier). Supposons
aussi \deleted{\illegible{1}} que \deleted{$\underline{\mathrm{Pic}}^0_{X_0/k}$ soit lisse, soit que}
les $\underline{\mathrm{Pic}}^0_{X_n/S_n}$ soient plats \deleted{et que} \inserted{(p. ex. $\underline{\mathrm{Pic}}^0_{X_0/k}$ lisse)} \unsure{de sorte|si bien} qu'on
peut former $\underline{NS}_{X_n/S_n}=\underline{\mathrm{Pic}}_{X_n/\unsure{k_n|S_n}}/\underline{\mathrm{Pic}}^0_{X_n/S_n}$, et
le sous-schéma $P_n=\underline{NS}_{X_n/S_n}^{+}$ (partie ample).
Les $P_n$ sont \deleted{\illegible{1}} \unsure{nets|étales|ouverts} sur les $S_n$, et par \unsure{là|la}
\unsure{compatibles|suite ils} définissent un schéma formel $\widehat P$ sur $\widehat S$,
donc que $\widehat P=\coprod_{\alpha}\widehat P_\alpha$
les $P_\alpha$ étant finis \inserted{nets} sur $S$, et connexes. Procédant
comme plus haut, on trouve que les conditions suivantes
sont équivalentes

(i) $\exists\ S'\to S$ \deleted{\unsure{propre}} \inserted{propre} surjectif, \inserted{\deleted{\illegible{1}}} avec $\widehat X_{\widehat S'}$ algébrique
et \uline{projectif} sur $S'$.

(ii) $\forall$ sous-schéma fermé \inserted{$T$ de $S$} de $\dim\leq 1$, si $T'$ est son normalisé,
alors $\widehat X_{\widehat T'}$ est algébrique et projectif sur $\widehat T'$.

(iii) $\exists$ morphisme \inserted{fini} surjectif net $S'\to S$, tel que
$\widehat X_{\widehat S'}$ soit algébrique et projectif.
% slanted
\margin{\uline{NB} Cette condition n'est bonne telle quelle que si $\underline{\mathrm{Pic}}^0_{\unsure{X_0/S_0|X_0/k}}$ est lisse; sinon \unsure{on sera} \deleted{\unsure{obligé}} obligé de prendre \inserted{(iii) et (iv)} $S'\to S$ \uline{fini et plat} (\unsure{car une polarisation|au sens f. loc. libre} \unsure{ne définit pas tjrs} un faisceau inversible\ldots}

(iv) Il existe un morphisme fini étale surjectif $S'\to S$, tel que pour tte composante irréductible
$S'_i$ de $S'$, avec str. réduite induite, on a $\widehat X_{\widehat S'_i}$ est
algébrique et projectif

\deleted{Ces conditions \unsure{impliquent la suivante}\unsure{.} \unsure{cette dernière}} De
plus, si $S'\to S$ est comme dans (iv), $\widehat X_{\widehat S'}$ est algébrique.
\deleted{\unsure{(iv$'$)} \unsure{(iv$'$)|(vi$'$)|(iv)} $\exists$ morphisme fini étale surjectif}
% slanted
\margin{\deleted{\unsure{Le}} \unsure{Les} deux dernières assertions}
% archive sheet number: 40
